\frfilename{mt133.tex}
\versiondate{29.3.10}
\copyrightdate{1994}

\def\chaptername{Pelengkap}
\def\sectionname{Konsep integrasi yang lebih luas}

\newsection{133}

Ada berbagai konteks yang membuat pemberian nilai kepada integral suatu
fungsi yang tidak sepenuhnya tercakup oleh definisi dasar dalam 122M menjadi
berguna. Dalam bagian ini saya menawarkan beberapa cara untuk menetapkan nilai
$\pm\infty$ kepada integral fungsi bernilai real (133A), mengintegralkan
fungsi bernilai kompleks (133C–133H), serta mendefinisikan integral atas dan
integral bawah (133I–133L). Dalam \S135 di bawah saya akan membahas
pengembangan lebih lanjut dari gagasan-gagasan Bab 12.

\leader{133A}{Integral tak hingga} Lazimnya frasa `$f$ terintegralkan'
dibatasi pada fungsi $f$ yang dapat diberi integral berhingga $\int f$
\cmmnt{ (sama seperti suatu deret disebut `dapat dijumlahkan' hanya jika
deret itu dapat diberi jumlah berhingga)}. Namun, untuk fungsi nonnegatif
kadang-kadang lebih mudah menuliskan `$\int f=\infty$' jika, dalam suatu
pengertian, satu-satunya alasan $f$ gagal terintegralkan adalah bahwa
integralnya terlalu besar; yaitu, $f$ didefinisikan hampir di mana-mana,
terukur secara virtual terhadap $\mu$, dan salah satu dari hal berikut berlaku:

\Centerline{$\{x:x\in\dom f,\,f(x)\ge\epsilon\}$}

\noindent memuat suatu himpunan berukuran tak hingga untuk suatu
$\epsilon>0$, atau

\Centerline{$\sup\{\int h:h\text{ sederhana},\,h\leae f\}=\infty$.}

\noindent\cmmnt{(Bandingkan 122J.) }Dengan kaidah ini,\cmmnt{ kita tetap
akan memperoleh}

\Centerline{$\int f_1+f_2=\int f_1+\int f_2$,
\quad$\int cf=c\int f$}

\noindent kapan pun $c\in\coint{0,\infty}$ dan $f_1$, $f_2$, $f$ merupakan
fungsi nonnegatif yang masing-masing memiliki $\int f_1$, $\int f_2$,
$\int f$ terdefinisi dalam $[0,\infty]$.

Karena itu kita dapat\cmmnt{ mengulangi definisi 122M dan} mengatakan bahwa

\Centerline{$\int f_1-f_2=\int f_1-\int f_2$}

\noindent kapan pun $f_1$, $f_2$ merupakan fungsi bernilai real sedemikian
sehingga $\int f_1$, $\int f_2$ terdefinisi dalam $[0,\infty]$ dan tidak
keduanya tak hingga\cmmnt{; syarat terakhir diberlakukan untuk menghindari
keharusan menghitung $\infty-\infty$}.

Kita tetap memiliki kaidah-kaidah

\Centerline{$\int f+g=\int f+\int g$,
\quad$\int(cf)=c\int f$,
\quad$\int|f|\ge|\int f|$}

\noindent setidaknya ketika ruas kanan dapat ditafsirkan, dengan mengizinkan
$0\cdot\infty=0$, tetapi tanpa mengizinkan tafsiran apa pun bagi
$\infty-\infty$; dan $\int f\le\int g$ kapan pun kedua integral terdefinisi
dan $f\leae g$. \cmmnt{(Namun, tentu saja sekarang mungkin terdapat
$f\le g$ dan $\int f=\int g=\pm\infty$ tanpa $f$ dan $g$ sama hampir di
mana-mana.)}

Tetapkan $f^+(x)=\max(f(x),0)$, $f^-(x)=\max(-f(x),0)$ untuk $x\in\dom f$;
maka

\Centerline{$\int f=\infty\,\iff\,\int f^+=\infty$ dan $f^-$
terintegralkan,}

\Centerline{$\int f=-\infty\,\iff\,f^+$ terintegralkan dan
$\int f^-=\infty$.}

\cmmnt{Untuk gagasan lebih lanjut dalam arah ini, lihat \S135 di bawah.}

\leader{133B}{Fungsi dengan nilai eksepsional} Juga lebih mudah jika kita
mengizinkan sebagai fungsi `terintegralkan' fungsi-fungsi $f$ yang sesekali
mengambil nilai bukan real\cmmnt{ -- biasanya ketika suatu rumus untuk
$f(x)$ mengizinkan nilai `$\infty$' berdasarkan suatu konvensi}. Untuk fungsi
seperti itu saya akan menuliskan $\int f=\int\tilde f$ jika
$\int\tilde f$ terdefinisi, dengan

\Centerline{$\dom\tilde f=\{x:x\in\dom f,\,f(x)\in\Bbb R\}$,
\quad$\tilde f(x)=f(x)$ untuk $x\in\dom\tilde f$.}

\cmmnt{\noindent Karena dalam konvensi ini saya tetap mengharuskan
$\tilde f$ didefinisikan hampir di mana-mana dalam $X$, himpunan
$\{x:x\in\dom f,\,f(x)\notin\Bbb R\}$ harus terabaikan.}

\cmmnt{
\leader{133C}{Fungsi bernilai kompleks} Seluruh teori fungsi terukur dan
terintegralkan yang dikembangkan sejauh ini dikhususkan bagi fungsi bernilai
real. Tidak diperlukan gagasan baru yang berarti untuk menangani fungsi
bernilai kompleks, tetapi barangkali sebaiknya saya menguraikan beberapa
rinciannya, karena ada banyak penerapan di mana fungsi bernilai kompleks
merupakan konteks kerja yang paling alami.
}%end of comment

\leader{133D}{Definisi (a)} Misalkan $X$ suatu himpunan dan $\Sigma$ suatu
aljabar-$\sigma$ subhimpunan-subhimpunan $X$. Jika $D\subseteq X$ dan
$f:D\to\Bbb C$ suatu fungsi, kita mengatakan bahwa $f$ {\bf terukur} jika
bagian real dan bagian imajinernya $\Real f$, $\Imag f$ terukur\cmmnt{ dalam
pengertian 121B–121C}.

\medskip

{\bf (b)} Misalkan $(X,\Sigma,\mu)$ suatu ruang ukur. Jika $f$ suatu fungsi
bernilai kompleks yang didefinisikan pada subhimpunan koterabaikan dari $X$,
kita mengatakan bahwa $f$ {\bf terintegralkan} jika bagian real dan bagian
imajinernya terintegralkan, dan kemudian

\Centerline{$\int f=\int\Real f+i\int\Imag f$.}

\medskip

{\bf (c)} Misalkan $(X,\Sigma,\mu)$ suatu ruang ukur, $H\in\Sigma$, dan $f$
suatu fungsi bernilai kompleks yang didefinisikan pada subhimpunan $X$. Maka
$\int_Hf$ adalah $\int(f\restr H)d\mu_H$ jika ini terdefinisi dalam
pengertian (b), dengan mengambil ukuran subruang $\mu_H$ sebagaimana dalam
131A–131B.

\ifnum\stylenumber=11\ifresultsonly\eject\fi\fi

\leader{133E}{Lema} (a) Jika $X$ suatu himpunan, $\Sigma$ suatu
aljabar-$\sigma$ subhimpunan-subhimpunan $X$, dan $f$ serta $g$ merupakan
fungsi terukur bernilai kompleks dengan domain $\dom f$, $\dom g\subseteq X$,
maka

\quad(i) $f+g:\dom f\cap\dom g\to\Bbb C$ terukur;

\quad(ii) $cf:\dom f\to\Bbb C$ terukur, untuk setiap $c\in\Bbb C$;

\quad(iii) $f\times g:\dom f\cap\dom g\to\Bbb C$ terukur;

\quad(iv) $f/g:\{x:x\in\dom f\cap\dom g,\,g(x)\ne 0\}\to\Bbb C$ terukur;

\quad(v) $|f|:\dom f\to\Bbb R$ terukur.

{(b)} Jika $\sequencen{f_n}$ suatu barisan fungsi terukur bernilai kompleks
yang didefinisikan pada subhimpunan-subhimpunan $X$, maka
$f=\lim_{n\to\infty}f_n$ terukur, jika kita mengambil $\dom f$ sebagai

$$\eqalign{\{x:x\in\bigcup_{n\in\Bbb N}&\bigcap_{m\ge n}\dom f_m,
  \,\lim_{n\to\infty}f_n(x)\text{ ada dalam }\Bbb C\}\cr
&=\dom(\lim_{n\to\infty}\Real f_n)
  \cap\dom(\lim_{n\to\infty}\Imag f_n).\cr}$$

\proof{{\bf (a)} Semuanya langsung diperoleh dari 121E, jika Anda menuliskan
rumus untuk bagian real dan bagian imajiner dari $f+g,\ldots,|f|$ dalam
hubungannya dengan bagian real dan bagian imajiner dari $f$ dan $g$.

\medskip

{\bf (b)} Gunakan 121Fa.
}%end of proof of 133E

\leader{133F}{Proposisi} Misalkan $(X,\Sigma,\mu)$ suatu ruang ukur.

(a) Jika $f$ dan $g$ merupakan fungsi terintegralkan bernilai kompleks yang
didefinisikan pada subhimpunan-subhimpunan koterabaikan dari $X$, maka $f+g$
dan $cf$ terintegralkan, $\int f+g=\int f+\int g$ dan
$\int cf=c\int f$, untuk setiap $c\in\Bbb C$.

{(b)} Jika $f$ suatu fungsi bernilai kompleks yang didefinisikan pada
subhimpunan koterabaikan dari $X$, maka $f$ terintegralkan jika dan hanya jika
$|f|$ terintegralkan dan $f$ terukur secara virtual terhadap $\mu$
\cmmnt{, yaitu, $\Real f$ dan $\Imag f$ terukur secara virtual terhadap
$\mu$}.

\proof{{\bf (a)} Gunakan 122Oa–122Ob.

\medskip

{\bf (b)} Pokoknya adalah bahwa $|\Real f|$,
$|\Imag f|\le|f|\le|\Real f|+|\Imag f|$; sekarang kita hanya perlu
menerapkan 122P sebanyak yang diperlukan.
}%end of proof of 133F

\leader{133G}{Teorema Konvergensi Terdominasi Lebesgue} Misalkan
$(X,\Sigma,\mu)$ suatu ruang ukur dan $\langle f_n\rangle_{n\in\Bbb N}$ suatu
barisan fungsi terintegralkan bernilai kompleks pada $X$ sedemikian sehingga
$f(x)=\lim_{n\to\infty}f_n(x)$ ada dalam $\Bbb C$ untuk hampir setiap
$x\in X$. Andaikan pula terdapat fungsi terintegralkan bernilai real $g$
pada $X$ sedemikian sehingga $|f_n|\leae g$ untuk setiap $n$. Maka $f$
terintegralkan dan $\lim_{n\to\infty}\int f_n$ ada serta sama dengan
$\int f$.

\proof{Terapkan 123C pada barisan-barisan $\sequencen{\Real f_n}$ dan
$\sequencen{\Imag f_n}$.
}%end of proof of 133G

\leader{133H}{Korolari} Misalkan $(X,\Sigma,\mu)$ suatu ruang ukur dan
$\ooint{a,b}$ suatu interval terbuka tak-kosong dalam $\Bbb R$. Misalkan
$f:X\times\ooint{a,b}\to\Bbb C$ suatu fungsi sedemikian sehingga

\inset{(i) integral $F(t)=\int f(x,t)dx$ terdefinisi untuk setiap
$t\in\ooint{a,b}$;}

\inset{(ii) turunan parsial $\pd{f}{t}$ dari $f$ terhadap variabel kedua
terdefinisi di setiap titik dalam $X\times\ooint{a,b}$;}

\inset{(iii) terdapat fungsi terintegralkan $g:X\to\coint{0,\infty}$
sedemikian sehingga $|\pd{f}{t}(x,t)|\le g(x)$ untuk setiap $x\in X$,
$t\in\ooint{a,b}$.}

\noindent Maka turunan $F'(t)$ dan integral
$\biggerint\pd{f}{t}(x,t)dx$ ada untuk setiap $t\in\ooint{a,b}$, dan
keduanya sama.

\proof{Terapkan 123D pada $\Real f$ dan $\Imag f$.}


\leader{133I}{Integral atas dan integral bawah}\dvArevised{2007}\cmmnt{
Sekarang saya kembali ke fungsi bernilai real.} Misalkan $(X,\Sigma,\mu)$
suatu ruang ukur dan $f$ suatu fungsi bernilai real yang didefinisikan hampir
di mana-mana dalam $X$. {\bf Integral atasnya} adalah

\Centerline{$\overlineint f
=\inf\{\int g:\int g$ terdefinisi dalam pengertian 133A dan $f\leae g\}$,}

\noindent dengan mengizinkan $\infty$ untuk $\inf\{\infty\}$ atau
$\inf\emptyset$ dan $-\infty$ untuk $\inf\Bbb R$. Demikian pula,
{\bf integral bawah} dari $f$ adalah

\Centerline{$\underlineint f
=\sup\{\int g:\int g$ terdefinisi, $f\geae g\}$,}

\noindent dengan mengizinkan $-\infty$ untuk $\sup\{-\infty\}$ atau
$\sup\emptyset$ dan $\infty$ untuk $\sup\Bbb R$.

\leader{133J}{Proposisi}\dvArevised{2007} Misalkan $(X,\Sigma,\mu)$ suatu
ruang ukur.

(a) Misalkan $f$ suatu fungsi bernilai real yang didefinisikan hampir di
mana-mana dalam $X$.

\quad (i) Jika $\overline{\int}f$ berhingga, terdapat fungsi terintegralkan
$g$ sedemikian sehingga $f\leae g$ dan $\int g=\overline{\int}f$. Dalam
hal ini,

\Centerline{$\{x:x\in\dom f\cap\dom g$, $g(x)\le f(x)+g_0(x)\}$}

\noindent berukuran luar penuh untuk setiap fungsi terukur
$g_0:X\to\ooint{0,\infty}$.

\quad (ii) Jika $\underline{\int}f$ berhingga, terdapat fungsi
terintegralkan $h$ sedemikian sehingga $h\leae f$ dan
$\int h=\underline{\int}f$. Dalam hal ini,

\Centerline{$\{x:x\in\dom f\cap\dom h$, $f(x)\le h(x)+h_0(x)\}$}

\noindent berukuran luar penuh untuk setiap fungsi terukur
$h_0:X\to\ooint{0,\infty}$.

{(b)} Untuk sembarang fungsi bernilai real $f$, $g$ yang didefinisikan pada
subhimpunan-subhimpunan koterabaikan dari $X$ dan sembarang $c\ge 0$,

\vskip 2pt

\quad(i) $\underline{\int}f\le\overline{\int}f$,

\vskip2pt

\quad (ii) $\overline{\int}f+g\le\overline{\int}f+\overline{\int}g$,

\vskip2pt

\quad (iii) $\overline{\int}cf=c\overline{\int}f$,

\vskip2pt

\quad (iv) $\underline{\int}(-f)=-\overline{\int}f$,

\vskip2pt

\quad (v) $\underline{\int}f+g\ge\underline{\int}f+\underline{\int}g$,

\vskip2pt

\quad (vi) $\underline{\int}cf=c\underline{\int}f$

\vskip2pt

\noindent kapan pun ruas kanan tidak melibatkan penjumlahan $\infty$ dengan
$-\infty$.

(c) Jika $f\leae g$ maka $\overline{\int}f\le\overline{\int}g$ dan
$\underline{\int}f\le\underline{\int}g$.

(d) Suatu fungsi bernilai real $f$ yang didefinisikan hampir di mana-mana
dalam $X$ terintegralkan jika dan hanya jika

\Centerline{$\overlineint f=\underlineint f=a\in\Bbb R$,}

\noindent dan dalam hal ini $\int f=a$.

(e) $\mu^*A=\overline{\int}\chi A$ untuk setiap $A\subseteq X$.

\proof{{\bf (a)(i)} Untuk setiap $n\in\Bbb N$, pilih fungsi $g_n$
sedemikian sehingga $f\leae g_n$ dan $\int g_n$ terdefinisi serta paling
besar $2^{-n}+\overline{\int}f$; karena
$\overline{\int}f\le\int g_n$, $\int g_n$ berhingga, sehingga $g_n$
terintegralkan. Tetapkan $h_n=\inf_{i\le n}g_i$ untuk setiap $n$; maka $h_n$
terintegralkan (karena
$|h_n-g_0|\le\sum_{i=0}^n|g_i-g_0|$ pada
$\bigcap_{i\le n}\dom g_i$), dan $f\leae h_n$, sehingga

\Centerline{$\overlineint f\le\int h_n
\le\int g_n\le 2^{-n}+\overlineint f$.}

\noindent Menurut teorema B. Levi (123A), yang diterapkan pada
$\sequencen{-h_n}$, $g(x)=\inf_{n\in\Bbb N}h_n(x)\in\Bbb R$ untuk hampir
setiap $x$, dan
$\int g=\inf_{n\in\Bbb N}\int h_n=\overline{\int}f$; juga, tentu saja,
$f\leae g$.

Sekarang ambil fungsi terukur $g_0:X\to\ooint{0,\infty}$, dan tinjau
himpunan

\Centerline{$A=\{x:x\in\dom f\cap\dom g$, $g(x)\le f(x)+g_0(x)\}$.}

\noindent\Quer\ Jika $A$ tidak berukuran luar penuh, terdapat himpunan
terukur tak-terabaikan $F\subseteq X\setminus A$. Karena $g_0$ positif tegas,
$F=\bigcup_{n\in\Bbb N}F_n$ dengan
$F_n=\{x:x\in F$, $g_0(x)\ge 2^{-n}\}$,
dan terdapat $n\in\Bbb N$ sedemikian sehingga $\mu F_n>0$. Tinjau fungsi
$g_1=g-2^{-n}\chi F_n$. Maka $f\leae g_1$. Selain itu,
$\int g_1=\int g-2^{-n}\mu F_n$ lebih kecil secara tegas daripada
$\int g$, sehingga $\overline{\int}f<\int g$.\ \Bang

\medskip

\quad{\bf (ii)} Gunakan argumen serupa, atau gunakan (b-iv).

\medskip

{\bf (b)(i)} Jika $\underline{\int}f=-\infty$ atau
$\overline{\int}f=\infty$, hal ini langsung. Jika tidak, hasilnya segera
mengikuti fakta bahwa apabila $g\leae f\leae h$, maka
$\int g\le\int h$ jika integral-integralnya terdefinisi (dalam arti luas).

\medskip

\quad{\bf (ii)} Jika $a>\overline{\int}f+\overline{\int}g$, baik
$\overline{\int}f$ maupun $\overline{\int}g$ tidak mungkin $\infty$, sehingga
harus ada fungsi $f_1$, $g_1$ sedemikian sehingga $f\leae f_1$,
$g\leae g_1$ dan $\int f_1+\int g_1\le a$. Sekarang
$f+g\leae f_1+g_1$, sehingga

\Centerline{$
\overlineint f+g\le\int f_1+g_1\le a$.}

\noindent Karena $a$ sembarang, kita memperoleh hasilnya.

\medskip

\quad{\bf (iii)($\pmb{\alpha}$)} Jika $c=0$, hal ini langsung. {\bf
($\pmb{\beta}$)} Jika $c>0$ dan $a>c\overline{\int}f$, harus terdapat
$f_1$ sedemikian sehingga $f\leae f_1$ dan $c\int f_1\le a$. Sekarang
$cf\leae cf_1$ dan $\int cf_1\le a$, sehingga $\overline{\int}cf\le a$.
Karena $a$ sembarang, $\overline{\int}cf\le c\overline{\int}f$.
{\bf($\pmb{\gamma}$)} Masih dengan mengandaikan $c>0$, kita juga memiliki

\Centerline{$c\overlineint f=c\overlineint c^{-1}cf
\le cc^{-1}\overlineint cf=\overlineint cf$,}

\noindent sehingga diperoleh kesamaan.

\medskip

\quad{\bf (iv)} Ini hanya karena $\int(-f_1)=-\int f_1$ bagi setiap fungsi
$f_1$ yang salah satu integralnya terdefinisi.

\medskip

\quad{\bf (v)-(vi)} Gunakan (iv) untuk mengubah $\underline{\int}$ menjadi
$\overline{\int}$, lalu terapkan (ii) atau (iii).

\medskip

{\bf (c)} Pernyataan-pernyataan ini langsung dari definisi, karena
(misalnya) jika $g\leae h$ maka $f\leae h$.

\medskip

{\bf (d)} Jika $f$ terintegralkan, maka

\Centerline{$\overlineint f=\int f=\underlineint f$}

\noindent menurut 122Od. Jika
$\overline{\int}f=\underline{\int}f=a\in\Bbb R$, maka, menurut (a), terdapat
fungsi terintegralkan $g$, $h$ sedemikian sehingga $g\leae f\leae h$ dan
$\int g=\int h=a$, sehingga $g\eae h$, menurut 122Rc, $g\eae f\eae h$,
dan $f$ terintegralkan, menurut 122Rb.

\medskip

{\bf (e)} Jika $E\supseteq A$ terukur, maka

\Centerline{$\mu E=\int\chi E\ge\overlineint\chi A$;}

\noindent karena $E$ sembarang, $\mu^*A\ge\overline{\int}\chi A$.
Jika $\int g$ terdefinisi dan $\chi A\leae g$, misalkan $E\subseteq\dom g$
suatu himpunan terukur koterabaikan sedemikian sehingga $g\restr E$ terukur,
dan tetapkan $F=\{x:x\in E$, $g(x)\ge 1\}$. Maka $A\setminus F$
terabaikan, sehingga $\mu^*A\le\mu F\le\int g$; karena $g$ sembarang,
$\mu^*A\le\overline{\int}\chi A$.
}%end of proof of 133J

\cmmnt{\medskip

\noindent{\bf Catatan} Saya harap rumus-rumus di sini mengingatkan Anda
pada $\limsup$, $\liminf$.}

\leader{133K}{Teorema konvergensi untuk integral
\dvrocolon{atas}}\cmmnt{ Kita memiliki versi-versi berikut dari teorema
B. Levi dan Lema Fatou.

\medskip

\noindent}{\bf Proposisi} Misalkan $(X,\Sigma,\mu)$ suatu ruang ukur, dan
$\sequencen{f_n}$ suatu barisan fungsi bernilai real yang didefinisikan
hampir di mana-mana dalam $X$.

(a) Jika, untuk setiap $n$, $f_n\leae f_{n+1}$, dan
$-\infty<\sup_{n\in\Bbb N}\overline{\int}f_n<\infty$, maka
$f(x)=\sup_{n\in\Bbb N}f_n(x)$ terdefinisi dalam $\Bbb R$ untuk hampir
setiap $x\in X$, dan
$\overline{\int}f=\sup_{n\in\Bbb N}\overline{\int}f_n$.

(b) Jika, untuk setiap $n$, $f_n\ge 0$ hampir di mana-mana, dan
$\liminf_{n\to\infty}\overline{\int}f_n<\infty$, maka
$f(x)=\liminf_{n\to\infty}f_n(x)$ terdefinisi dalam $\Bbb R$ untuk hampir
setiap $x\in X$, dan
$\overline{\int}f\le\liminf_{n\to\infty}\overline{\int}f_n$.

\proof{{\bf (a)} Tetapkan
$c=\sup_{n\in\Bbb N}\overline{\int}f_n$. Untuk setiap $n$, terdapat fungsi
terintegralkan $g_n$ sedemikian sehingga $f_n\leae g_n$ dan
$\int g_n=\overline{\int}f_n$ (133J(a-i)).
Tetapkan $g'_n=\min(g_n,g_{n+1})$; maka $g'_n$ terintegralkan dan
$f_n\leae g'_n\leae g_n$, sehingga

\Centerline{$\overlineint f_n\le\int g'_n
\le\int g_n=\overlineint f_n$}

\noindent dan $g'_n$ harus sama dengan $g_n$ hampir di mana-mana.
Akibatnya $g_n\leae g_{n+1}$ untuk setiap $n$, sementara
$\sup_{n\in\Bbb N}\int g_n=c<\infty$.
Menurut teorema B. Levi, $g=\sup_{n\in\Bbb N}g_n$ terdefinisi sebagai
fungsi bernilai real hampir di mana-mana dalam $X$, dan $\int g=c$. Sekarang
tentu saja $f(x)$ terdefinisi, dan tidak lebih besar daripada $g(x)$, untuk
setiap $x\in\dom g\cap\bigcap_{n\in\Bbb N}\dom f_n$ sedemikian sehingga
$f_n(x)\le g_n(x)$ untuk setiap $n$, yaitu untuk hampir setiap $x$; sehingga
$\overline{\int}f\le\int g=c$. Di sisi lain, $f_n\leae f$, sehingga
$\overline{\int}f_n\le\overline{\int}f$ untuk setiap $n\in\Bbb N$;
akibatnya $\overline{\int}f$ sekurang-kurangnya $c$, dan karena itu sama
dengan $c$, sebagaimana diperlukan.

\medskip

{\bf (b)} Argumennya mengikuti argumen 123B. Tetapkan
$c=\liminf_{n\to\infty}\overline{\int}f_n$. Untuk setiap $n$, tetapkan
$g_n=\inf_{m\ge n}f_m$; maka
$\overline{\int}g_n\le\inf_{m\ge n}\overline{\int}f_m\le c$.
Kita memiliki $g_n(x)\le g_{n+1}(x)$ untuk setiap $x\in\dom g_n$, yaitu
hampir di mana-mana, untuk setiap $n$; sehingga, menurut (a),

\Centerline{$\overlineint g
=\sup_{n\in\Bbb N}\overlineint g_n\le c$,}

\noindent dengan

\Centerline{$g=\sup_{n\in\Bbb N}g_n\eae\liminf_{n\to\infty}f_n$,}

\noindent dan $\overline{\int}\liminf_{n\to\infty}f_n\le c$, sebagaimana
diklaim.
}%end of proof of 133K

\leader{*133L}{}\dvAnew{2007}\cmmnt{ Hasil berikut berada pada tingkat yang
kurang mendasar daripada hasil-hasil dalam 133J, tetapi tetap penting.

\medskip

\noindent}{\bf Proposisi}
Misalkan $(X,\Sigma,\mu)$ suatu ruang ukur dan $f$ suatu fungsi bernilai real
yang didefinisikan hampir di mana-mana dalam $X$. Andaikan $h_1$, $h_2$
merupakan fungsi nonnegatif yang terukur secara virtual dan didefinisikan
hampir di mana-mana dalam $X$. Maka

\Centerline{$\overlineint f\times(h_1+h_2)
=\overlineint f\times h_1+\overlineint f\times h_2$,}

\noindent dengan catatan bahwa di sini, untuk kali ini saja, kita dapat
menafsirkan $\infty+(-\infty)$ atau $(-\infty)+\infty$ sebagai $\infty$
jika hal itu diperlukan pada ruas kanan.

\proof{{\bf (a)} Jika $\overline{\int}f\times h_1=\infty$ atau
$\overline{\int}f\times h_2=\infty$, maka
$\overline{\int}f\times(h_1+h_2)=\infty$. \Prf\Quer\ Jika tidak, terdapat
$g$ sedemikian sehingga $f\times(h_1+h_2)\leae g$ dan $\int g<\infty$.
Dalam hal ini,

\Centerline{$f\times h_1
\leae f^+\times h_1
\leae f^+\times(h_1+h_2)
=(f\times(h_1+h_2))^+
\leae g^+$}

\noindent sehingga $\overline{\int}f\times h_1\le\int g^+<\infty$.
Demikian pula, $\overline{\int}f\times h_2<\infty$; bertentangan dengan
hipotesis kita.\ \Bang\QeD\ Jadi dalam kasus ini, berdasarkan kaidah lokal
$\infty+(-\infty)=(-\infty)+\infty=\infty$, kita memperoleh hasilnya.

\medskip

{\bf (b)} Sekarang andaikan integral atas $\overline{\int}f\times h_1$ dan
$\overline{\int}f\times h_2$ keduanya lebih kecil daripada $\infty$,
sehingga jumlah keduanya dapat ditafsirkan menurut kaidah biasa.
Menurut 133J(b-ii),
$\overline{\int}f\times(h_1+h_2)
\le\overline{\int}f\times h_1+\overline{\int}f\times h_2<\infty$.
Untuk arah sebaliknya, andaikan bahwa $g\geae f\times(h_1+h_2)$ dan
$\int g<\infty$. Untuk $i=1$, $2$, tetapkan

$$\eqalign{g_i(x)
&=\Bover{g(x)h_i(x)}{h_1(x)+h_2(x)}
  \text{ jika }x\in\dom g\cap\dom h_1\cap\dom h_2
  \text{ dan }h_1(x)+h_2(x)>0,\cr
&=0\text{ untuk titik-titik lain }x\in X.\cr}$$

\noindent Maka, untuk kedua $i$, $g_i$ terukur secara virtual,
$g_i^+\leae g^+$ dan $g_i\geae f\times h_i$; sementara $g\geae g_1+g_2$.
\Prf\ Himpunan

\Centerline{$H=\{x:x\in\dom f\cap\dom g\cap\dom h_1\cap\dom h_2$,
$g(x)\ge f(x)(h_1(x)+h_2(x))\}$}

\noindent bersifat koterabaikan, dan untuk $x\in H$

$$\eqalign{g(x)
&=g_1(x)+g_2(x)\text{ jika }h_1(x)+h_2(x)>0,\cr
&\ge 0=g_1(x)+g_2(x)\text{ jika }h_1(x)+h_2(x)=0.  \text{ \Qed}\cr}$$

\noindent Jadi

\Centerline{$\overlineint f\times h_1+\overlineint f\times h_2
\le\int g_1+\int g_2
=\int g_1+g_2
\le\int g$}

\noindent (karena $\int g_1$ dan $\int g_2$ keduanya paling besar
$\int g^+<\infty$, sehingga kita dapat menjumlahkannya menurut kaidah biasa).
Karena $g$ sembarang,
$\overline{\int}f\times h_1+\overline{\int}f\times h_2
\le\overline{\int}f\times(h_1+h_2)$ dan harus berlaku kesamaan.
}%end of proof of 133L

\exercises{
\leader{133X}{Latihan dasar $\pmb{>}$(a)}
%\spheader 133Xa
Misalkan $(X,\Sigma,\mu)$ suatu ruang ukur, dan
$f:X\to\coint{0,\infty}$ suatu fungsi terukur. Tunjukkan bahwa

$$\eqalign{\int fd\mu
&=\sup_{n\in\Bbb N}2^{-n}\sum_{k=1}^{4^n}\mu\{x:f(x)\ge 2^{-n}k\}\cr
&=\lim_{n\to\infty}2^{-n}\sum_{k=1}^{4^n}\mu\{x:f(x)\ge 2^{-n}k\}\cr}$$

\noindent dalam $[0,\infty]$.
%133A

\spheader 133Xb
Misalkan $(X,\Sigma,\mu)$ suatu ruang ukur dan $f$ suatu fungsi bernilai
kompleks yang didefinisikan pada subhimpunan $X$. (i) Tunjukkan bahwa jika
$E\in\Sigma$, maka $f\restr E$ terintegralkan terhadap $\mu_E$ jika dan hanya
jika $\tilde f$ terintegralkan terhadap $\mu$, dengan menuliskan $\mu_E$
untuk ukuran subruang pada $E$ dan $\tilde f(x)=f(x)$ jika
$x\in E\cap\dom f$, $0$ jika $x\in X\setminus E$; dan dalam hal ini
$\int_Efd\mu_E=\int \tilde fd\mu$. (ii) Tunjukkan bahwa jika $E\in\Sigma$
dan $f$ didefinisikan hampir di mana-mana terhadap $\mu$, maka
$f\restr E$ terintegralkan terhadap $\mu_E$ jika dan hanya jika
$f\times\chi E$ terintegralkan terhadap $\mu$, dan dalam hal ini
$\int_Ef=\int f\times\chi E$. (iii) Tunjukkan bahwa jika $\int_Ef=0$ untuk
setiap $E\in\Sigma$, maka $f=0$ hampir di mana-mana.
%133E

\spheader 133Xc Andaikan $(X,\Sigma,\mu)$ suatu ruang ukur dan $G$ suatu
subhimpunan terbuka dari $\Bbb C$, yaitu himpunan sedemikian sehingga untuk
setiap $w\in G$ terdapat $\delta>0$ sedemikian sehingga
$\{z:|z-w|<\delta\}\subseteq G$. Misalkan $f:X\times G\to\Bbb C$ suatu
fungsi, dan andaikan bahwa turunan $\pd{f}{z}$ dari $f$ terhadap variabel
kedua ada untuk semua $x\in X$, $z\in G$. Andaikan pula bahwa
(i) $F(z)=\int f(x,z)dx$ ada untuk setiap $z\in G$ (ii) terdapat fungsi
terintegralkan $g$ sedemikian sehingga $|\pd{f}{z}(x,z)|\le g(x)$ untuk
setiap $x\in X$, $z\in G$. Tunjukkan bahwa turunan $F'$ dari $F$ ada di
setiap titik dalam $G$, dan
$F'(z)=\biggerint\pd{f}{z}(x,z)dx$ untuk setiap $z\in G$.
({\it Petunjuk\/}: Anda perlu memeriksa bahwa
$|f(x,z)-f(x,w)|\le|z-w|g(x)$ kapan pun $x\in X$, $z\in G$, dan $w$ dekat
dengan $z$.)
%133F

\sqheader 133Xd Misalkan $f$ suatu fungsi bernilai kompleks yang
didefinisikan hampir di mana-mana pada $\coint{0,\infty}$, yang seperti biasa
dilengkapi dengan ukuran Lebesgue. {\bf Transformasi Laplace}-nya adalah
fungsi $F$ yang didefinisikan dengan menuliskan

\Centerline{$F(s)=\int_0^{\infty}e^{-sx}f(x)dx$}

\noindent untuk semua bilangan kompleks $s$ yang integralnya terdefinisi
dalam $\Bbb C$.

\quad(i) Tunjukkan bahwa jika $s\in\dom F$ dan $\Real s'\ge\Real s$, maka
$s'\in\dom F$ (karena $|e^{-s'x}e^{sx}|\le 1$ untuk semua $x$).

\quad(ii) Tunjukkan bahwa $F$ analitik (yaitu, terdiferensialkan sebagai
fungsi variabel kompleks) pada interior domainnya. \Hint{133Xc.}

\quad(iii) Tunjukkan bahwa jika $F$ terdefinisi di suatu titik, maka
$\lim_{\Real s\to\infty}F(s)=0$.

\quad(iv) Tunjukkan bahwa jika $f$, $g$ memiliki transformasi Laplace $F$,
$G$, maka transformasi Laplace dari $f+g$ adalah $F+G$, setidaknya pada
$\dom F\cap \dom G$.
%133Xc 133F

\sqheader 133Xe Misalkan $f$ suatu fungsi terintegralkan bernilai kompleks
yang didefinisikan hampir di mana-mana dalam $\Bbb R$, yang seperti biasa
dilengkapi dengan ukuran Lebesgue. {\bf Transformasi Fourier}-nya adalah
fungsi $\varhatf$ yang didefinisikan dengan

\Centerline{$\varhatf(s)=\Bover{1}{\sqrt{2\pi}}
\int_{-\infty}^{\infty}e^{-isx}f(x)dx$}

\noindent untuk semua $s$ real.

\quad(i) Tunjukkan bahwa $\varhatf$ kontinu. \Hint{gunakan Teorema
Konvergensi Terdominasi Lebesgue pada barisan-barisan berbentuk
$f_n(x)=e^{-is_nx}f(x)$.}

\quad(ii) Tunjukkan bahwa jika $f$, $g$ memiliki transformasi Fourier
$\varhatf$, $\varhat g$, maka transformasi Fourier dari $f+g$ adalah
$\varhatf+\varhat g$.

\quad(iii) Tunjukkan bahwa jika $\int xf(x)dx$ ada, maka $\varhatf$
terdiferensialkan, dengan
$\varhatf\vthsp'(s)=-\Bover{i}{\sqrt{2\pi}}\int xe^{-isx}f(x)dx$ untuk
setiap $s$.
%133Xc 133F

\spheader 133Xf Misalkan $(X,\Sigma,\mu)$ suatu ruang ukur dan
$\sequencen{f_n}$ suatu barisan fungsi bernilai real yang masing-masing
didefinisikan hampir di mana-mana dalam $X$. Andaikan terdapat fungsi
terintegralkan bernilai real $g$ sedemikian sehingga $|f_n|\leae g$ untuk
setiap $n$. Tunjukkan bahwa

\Centerline{$\overlineint\liminf_{n\to\infty}f_n
\le\liminf_{n\to\infty}\overlineint f_n$,\quad
$\underlineint\limsup_{n\to\infty}f_n
\ge\limsup_{n\to\infty}\underlineint f_n$.}
%133K


\leader{133Y}{Latihan lanjutan (a)}
%\spheader 133Ya
Gunakan gagasan-gagasan 133C–133H untuk mengembangkan teori fungsi terukur
dan terintegralkan yang mengambil nilai dalam $\BbbR^r$, dengan $r\ge 2$.
%133H

\spheader 133Yb Misalkan $X$ suatu himpunan dan $\Sigma$ suatu
aljabar-$\sigma$ subhimpunan-subhimpunan $X$. Misalkan $Y$ suatu subhimpunan
$X$ dan $f:Y\to\Bbb C$ suatu fungsi $\Sigma_Y$-terukur, dengan
$\Sigma_Y=\{E\cap Y:E\in\Sigma\}$. Tunjukkan bahwa terdapat fungsi
$\Sigma$-terukur $\tilde f:X\to\Bbb C$ yang memperluas $f$. \Hint{121I.}
%133F

\spheader 133Yc Misalkan $f$ suatu fungsi terintegralkan bernilai kompleks
yang didefinisikan hampir di mana-mana dalam $\BbbR^r$, yang seperti biasa
dilengkapi dengan ukuran Lebesgue, dengan $r\ge 1$. {\bf Transformasi
Fourier}-nya adalah fungsi $\varhatf$ yang didefinisikan dengan

\Centerline{$\varhatf(s)=\Bover{1}{{(\sqrt{2\pi})}^r}
\int e^{-is\dotproduct x}f(x)dx$}

\noindent untuk semua $s\in\BbbR^r$, dengan menuliskan $s\dotproduct x$
untuk $\sigma_1\xi_1+\ldots+\sigma_r\xi_r$ jika
$s=(\sigma_1,\ldots,\sigma_r)$,
$x=(\xi_1,\ldots,\xi_r)\in\BbbR^r$.

\quad(i) Tunjukkan bahwa $\varhatf$ kontinu.

\quad(ii) Tunjukkan bahwa jika $f$, $g$ memiliki transformasi Fourier
$\varhatf$, $\varhat g$, maka transformasi Fourier dari $f+g$ adalah
$\varhatf+\varhat g$.

\quad (iii) Tunjukkan bahwa jika $\int\|x\||f(x)|dx$ berhingga
(dengan mengambil $\|x\|=\sqrt{\xi_1^2+\ldots+\xi_r^2}$ jika
$x=(\xi_1,\dots,\xi_r)$), maka $\varhatf$ terdiferensialkan, dengan

\Centerline{$\Pd{\varhatf}{\sigma_k}(s)
=-\Bover{i}{{(\sqrt{2\pi})}^r}\int\xi_ke^{-is\dotproduct x}f(x)dx$}

\noindent untuk setiap $s\in\BbbR^r$, $k\le r$.
%133Xe 133F

\spheader 133Yd Ingat kembali definisi fungsi `kuasi-sederhana' dari 122Yd.
Tunjukkan bahwa untuk sembarang ruang ukur $(X,\Sigma,\mu)$ dan sembarang
fungsi bernilai real $f$ yang didefinisikan hampir di mana-mana dalam $X$,

\Centerline{$
\overlineint f=\inf\{\int g:g$ kuasi-sederhana, $f\leae g\}$,}

\Centerline{$
\underlineint f=\sup\{\int g:g$ kuasi-sederhana, $f\geae g\}$,}

\noindent dengan mengizinkan $\infty$ untuk $\inf\emptyset$ dan
$\sup\Bbb R$, serta $-\infty$ untuk $\inf\Bbb R$ dan $\sup\emptyset$.
%133J

\spheader 133Ye Nyatakan dan buktikan hasil serupa mengenai fungsi-fungsi
`pseudo-sederhana' dari 122Ye.
%133Yd 133J
}%end of exercises

\endnotes{
\Notesheader{133} Saya telah menguraikan bagian ini secara terperinci,
meskipun di dalamnya tidak ada sesuatu yang sungguh-sungguh dapat disebut
gagasan baru, karena bagian ini memberi kita kesempatan untuk meninjau kembali
pekerjaan sebelumnya, dan karena manipulasi-manipulasi yang sekarang, saya
harap, mulai tampak `jelas' bagi Anda sebenarnya hanya dapat dibenarkan
melalui teorema-teorema yang sulit; saya percaya bahwa setidaknya sesekali
kita patut melihat kembali titik-titik tepat tempat pembenaran itu dituliskan.

Anda mungkin telah memperhatikan kemiripan antara hasil-hasil yang melibatkan
`integral atas', sebagaimana diuraikan di sini, dan hasil-hasil dalam \S132
mengenai `ukuran luar' (misalnya 132Ae dan 133Ka, atau 132Xe dan 133Kb).
Kemiripan ini bukan kebetulan; semacam penjelasan dapat ditemukan dalam 252Ym
di Jilid 2.
}%end of notes

\discrpage
